\NativeSlot{47}{36}{%
% p018_s01

\ExamQuestion{8（续）}{7}%
新疆棉毛巾，一共买了8盒，每盒有3条，每盒毛巾单价是\mbox{（\hspace{1.8em}）}元，平均每条毛巾\mbox{（\hspace{1.8em}）}元。\par
}
\NativeSlot{98}{73}{%
% p018_s02

\ExamQuestion{9.}{5}%
国庆出游，王叔叔开车带家人自驾游，其中高速路段车速平均每小时106千米，全程累计开了22小时，一共开了多少千米？\par\noindent （1）解决这个问题所依据的数量关系式是\mbox{（\hspace{12em}）}。\par

% p018_s03

\ExamQuestion{（2）}{5}%
右图竖式中，第②步计算是\(106\times\)\mbox{（\hspace{1.8em}）}，表示的是\mbox{（\hspace{1.8em}）}小时行驶的路程。\par
\BookFigure{figures/p018_s03-01.tex}
}
\NativeSlot{225}{75}{%
\ExamGroup{二、选择（将正确答案的序号填在括号里）。{\normalfont\kaishu\fontsize{10.8bp}{14bp}\selectfont （每空2分，共14分）}}

% p018_s04

\ExamQuestion{1.}{5}%
小华步行去商店，路上用了10分钟；买了3本练习本，花了15元。这里的10分钟、15元、3本分别表示的是\mbox{（\hspace{1.8em}）}。\par
\ExamOptions{A. 路程，总价，数量}{B. 时间，总价，单价}{C. 时间，总价，数量}{D. 速度，单价，数量}
}
\NativeSlot{350}{45}{%
% p018_s05

\ExamQuestion{2.}{5}%
每瓶矿泉水3元，买12瓶一共花了多少钱？用\mbox{（\hspace{1.8em}）}计算；一名外卖骑手，5分钟骑行1400米，求速度，用\mbox{（\hspace{1.8em}）}计算。\par
\ExamOptions{A. 加法}{B. 乘法}{C. 除法}{D. 减法}
}
\NativeSlot{426}{59}{%
% p018_s06

\ExamQuestion{3.}{5}%
如图，一列货车从甲地开往乙地用了4个小时，货车的速度是\mbox{（\hspace{1.8em}）}。\par
\ExamOptions{A. 70时}{B. 70千米}{C. 70千米/时}{D. 280千米/时}
\BookFigure{figures/p018_s06-01.tex}
}
\NativeSlot{526}{44}{%
% p018_s07

\ExamQuestion{4.}{5}%
下面问题中，适合用“速度\(\times\)时间”来解决的是\mbox{（\hspace{1.8em}）}。\par
\ExamOptions{A. 求每支钢笔的钱数}{B. 求汽车5小时行驶的千米数}{C. 求小明每分钟步行的米数}{D. 计算长方形的周长}
}
\NativeSlot{601}{116}{%
% p018_s08

\ExamQuestion{5.}{5}%
这个公路指示牌告诉我们当前位置距青岛的\mbox{（\hspace{1.8em}）}还有203千米。\par
\ExamOptions{A. 时间}{B. 速度}{C. 路程}{D. 数量}
\par\smallskip\Needspace{4\baselineskip}\begingroup\small\setlength{\tabcolsep}{3pt}
\noindent\begin{tabularx}{\linewidth}{|*{2}{>{\centering\arraybackslash}X|}}\hline
目的地 & 距离\\\hline
青岛 & 203千米\\\hline
\end{tabularx}\par
\endgroup\smallskip

% p018_s09

\ExamQuestion{6.}{5}%
张平用18元买了6本笔记本，小华买这样的4本笔记本花了多少钱？列式正确的是\mbox{（\hspace{1.8em}）}。\par
\ExamOptions{A. \(18\times4\div6\)}{B. \(18\times6\div4\)}{C. \(18\div6\times4\)}{D. \(18\div6+4\)}
}
